% GENERATED FILE. Edit canonical lessons, then run npm run generate:books.
% canonical-source-hash: fnv1a64:f704d6664db02ef8
% canonical-lessons: TE-C84-fever, TE-C84-pain, TE-C84-tiredness, TE-C84-thirst, TE-C84-rest

\chapter{When You Are Not Well}
\label{ch:te-not-well}
\hlchaptermodality{84}

\begin{chapteropening}
\textbf{By the end of this chapter:} \emph{I can say that I have a fever, that something hurts, that I am tired or thirsty, and ask someone to rest.}

Complete TE-C84-rest: name what is wrong and tell a tired or feverish friend to rest.
\end{chapteropening}

\section[jvaraṁ]{\te{జ్వరం} (jvaraṁ) — a fever}

\label{lesson:TE-C84-fever}



\begin{quote}
Before the new one: say the Telugu for a disease, then the Telugu for health.
\end{quote}

\subsection*{You'll want to know: \te{జ్వరం}}
\textbf{\te{జ్వరం}} (\emph{jvaraṁ}) — \textquotedblleft{}a fever\textquotedblright{}.

\textbf{\te{నాకు} \te{జ్వరం} \te{వచ్చింది}} — \textquotedblleft{}I have a fever\textquotedblright{}, literally \textquotedblleft{}to me a fever has come\textquotedblright{}.

Look at the shape. The person sits in the dative \te{నాకు} and the fever \emph{arrives}, on the same come-verb as \textbf{\te{నాకు} \te{తెలుగు} \te{వచ్చు}}. In Telugu an illness is something that comes to you.

\begin{grammarlens}[title={What you've built}]
Your first word for what is wrong, on the pattern you already use for knowing Telugu.
\end{grammarlens}

\subsection*{Your turn}
\begin{itemize}
\raggedright
  \item \emph{Say it:} \emph{jvaraṁ}
  \item \emph{Say it:} \emph{jvaraṁ}, once more
  \item \emph{Say it:} say \emph{nāku jvaraṁ vaccindi}, I have a fever
  \item \emph{Recall:} say the Telugu for a disease, then the Telugu for health, then say \emph{jvaraṁ} again
\end{itemize}

\begin{tcolorbox}[breakable,colback=teal!4,colframe=teal!35!black,title={Before you move on}]
What does \te{జ్వరం} mean? (\textquotedblleft{}A fever\textquotedblright{}.) Say it once more.
\end{tcolorbox}
\section[noppi]{\te{నొప్పి} (noppi) — a pain, an ache}

\label{lesson:TE-C84-pain}



\begin{quote}
Before the new one: say the Telugu for health, then the Telugu for a fever.
\end{quote}

\subsection*{You'll want to know: \te{నొప్పి}}
\textbf{\te{నొప్పి}} (\emph{noppi}) — \textquotedblleft{}a pain, an ache\textquotedblright{}.

Put a body part in front of it and you have the complaint: \textbf{\te{తల} \te{నొప్పి}} (\emph{tala noppi}), a headache, and \textbf{\te{కడుపు} \te{నొప్పి}} (\emph{kaḍupu noppi}), a stomach ache. Both body words are ones you already own.

\begin{grammarlens}[title={What you've built}]
One word that turns every body part you know into something you can complain about.
\end{grammarlens}

\subsection*{Your turn}
\begin{itemize}
\raggedright
  \item \emph{Say it:} \emph{noppi}
  \item \emph{Say it:} \emph{noppi}, once more
  \item \emph{Say it:} say \emph{tala noppi}, then \emph{kaḍupu noppi}
  \item \emph{Recall:} say the Telugu for health, then the Telugu for a fever, then say \emph{noppi} again
\end{itemize}

\begin{tcolorbox}[breakable,colback=teal!4,colframe=teal!35!black,title={Before you move on}]
What does \te{నొప్పి} mean? (\textquotedblleft{}A pain, an ache\textquotedblright{}.) Say it once more.
\end{tcolorbox}
\section[alasaṭa]{\te{అలసట} (alasaṭa) — tiredness}

\label{lesson:TE-C84-tiredness}



\begin{quote}
Before the new one: say the Telugu for a fever, then the Telugu for a pain.
\end{quote}

\subsection*{You'll want to know: \te{అలసట}}
\textbf{\te{అలసట}} (\emph{alasaṭa}) — \textquotedblleft{}tiredness\textquotedblright{}.

\textbf{\te{నాకు} \te{అలసటగా} \te{ఉంది}} — \textquotedblleft{}I feel tired\textquotedblright{}, literally \textquotedblleft{}to me it is tired-ly\textquotedblright{}.

The ending \textbf{-\te{గా}} turns the noun into a manner, and \textbf{\te{ఉంది}} says \textquotedblleft{}it is\textquotedblright{}. This frame, \emph{\te{నాకు} ... \te{గా} \te{ఉంది}}, is how Telugu says many of the things English says with \emph{I am}.

\begin{grammarlens}[title={What you've built}]
The frame \te{నాకు} ... \te{గా} \te{ఉంది}, which the next words keep using.
\end{grammarlens}

\subsection*{Your turn}
\begin{itemize}
\raggedright
  \item \emph{Say it:} \emph{alasaṭa}
  \item \emph{Say it:} \emph{alasaṭa}, once more
  \item \emph{Say it:} say \emph{nāku alasaṭagā undi}, I feel tired
  \item \emph{Recall:} say the Telugu for a fever, then the Telugu for a pain, then say \emph{alasaṭa} again
\end{itemize}

\begin{tcolorbox}[breakable,colback=teal!4,colframe=teal!35!black,title={Before you move on}]
What does \te{అలసట} mean? (\textquotedblleft{}Tiredness\textquotedblright{}.) Say it once more.
\end{tcolorbox}
\section[dāhaṁ]{\te{దాహం} (dāhaṁ) — thirst}

\label{lesson:TE-C84-thirst}



\begin{quote}
Before the new one: say the Telugu for a pain, then the Telugu for tiredness.
\end{quote}

\subsection*{You'll want to know: \te{దాహం}}
\textbf{\te{దాహం}} (\emph{dāhaṁ}) — \textquotedblleft{}thirst\textquotedblright{}.

\textbf{\te{నాకు} \te{దాహంగా} \te{ఉంది}} — \textquotedblleft{}I am thirsty\textquotedblright{}. Same frame as tiredness, a new noun in the middle.

And you can already ask for the cure: \textbf{\te{నాకు} \te{నీళ్ళు} \te{కావాలి}}, \textquotedblleft{}I want water\textquotedblright{}.

\begin{grammarlens}[title={What you've built}]
Thirst, and the request that answers it.
\end{grammarlens}

\subsection*{Your turn}
\begin{itemize}
\raggedright
  \item \emph{Say it:} \emph{dāhaṁ}
  \item \emph{Say it:} \emph{dāhaṁ}, once more
  \item \emph{Say it:} say \emph{nāku dāhaṁgā undi}, then ask for water
  \item \emph{Recall:} say the Telugu for a pain, then the Telugu for tiredness, then say \emph{dāhaṁ} again
\end{itemize}

\begin{tcolorbox}[breakable,colback=teal!4,colframe=teal!35!black,title={Before you move on}]
What does \te{దాహం} mean? (\textquotedblleft{}Thirst\textquotedblright{}.) Say it once more.
\end{tcolorbox}
\section[viśrānti]{\te{విశ్రాంతి} (viśrānti) — rest}

\label{lesson:TE-C84-rest}



\begin{quote}
Before the new one: say the Telugu for tiredness, then the Telugu for thirst.
\end{quote}

\subsection*{You'll want to know: \te{విశ్రాంతి}}
\textbf{\te{విశ్రాంతి}} (\emph{viśrānti}) — \textquotedblleft{}rest\textquotedblright{}.

\textbf{\te{విశ్రాంతి} \te{తీసుకోండి}} — \textquotedblleft{}please take some rest\textquotedblright{}, what you say to someone with a fever or a long day behind them. The polite ending \textbf{-\te{అండి}} is the one you already put on \emph{raṇḍi}.

\begin{grammarlens}[title={What you've built}]
Something to say to the person who is unwell, not only about yourself.
\end{grammarlens}

\subsection*{Your turn}
\begin{itemize}
\raggedright
  \item \emph{Say it:} \emph{viśrānti}
  \item \emph{Say it:} \emph{viśrānti}, once more
  \item \emph{Say it:} say \emph{viśrānti tīsukōṇḍi} to a tired friend
  \item \emph{Recall:} say the Telugu for tiredness, then the Telugu for thirst, then say \emph{viśrānti} again
\end{itemize}

\begin{tcolorbox}[breakable,colback=teal!4,colframe=teal!35!black,title={Before you move on}]
What does \te{విశ్రాంతి} mean? (\textquotedblleft{}Rest\textquotedblright{}.) Say it once more.
\end{tcolorbox}
